\documentclass[11pt]{amsart}% src: user assignment: document-class requirement
\usepackage{amsmath}
\usepackage{amssymb}
\usepackage{amsthm}
\usepackage{geometry}
\usepackage{hyperref}
\usepackage{booktabs}
\geometry{margin=1in}
\hypersetup{hidelinks}

\newcommand{\Q}{\mathbb{Q}}
\newcommand{\Z}{\mathbb{Z}}
\newcommand{\F}{\mathbb{F}}
\DeclareMathOperator{\supp}{supp}
\newcommand{\leg}[2]{\left(\frac{#1}{#2}\right)}
\newtheorem{proposition}{Proposition}
\newtheorem{remark}{Remark}

\title{A Machine-Verified Finite Program:\\
Engine Discipline, Artifact Schemas, and Reproducibility Scope}
\author{The H10Q verification program}
\date{}
% src: README.md:1

\begin{document}
\begin{abstract}
This companion describes what the finite H10Q computation verifies, how it
refuses arithmetic it cannot certify, and how its persisted artifacts may be
replayed.  The arithmetic kernel admits primality only through a deterministic
Miller--Rabin range or a Pocklington certificate, including the
Brillhart--Lehmer--Selfridge relaxation.  The surrounding suite separates
machine proofs from conditional deductions and measured evidence, freezes its
finite authorities in code, and records enough schema information to prevent a
search log from silently becoming a theorem.  We give the exact downstream
regeneration commands, the licenses attached to each artifact row, and the
limitations left by missing historical producers and finite prime windows.
% src: README.md:12-43; README.md:191-221
\end{abstract}

\section{The verification contract}

The program is a verifier of finite mathematical statements, not an oracle for
Hilbert's Tenth Problem over \(\Q\).  A positive row is accepted only after its
stated hypotheses and exact identities have been replayed.  A bounded scan
licenses a statement about the scanned rows and no statement about the
unscanned complement.  In particular, the canonical grid is closed
instance-by-instance, while the uniform all-\(w\) assertion remains an analytic
frontier; the audited simultaneous prime-value input has degree \(8\) and is
still open.
% src: README.md:138-145; README.md:225-231

\subsection{The label calculus}
The ledger uses three labels as a small calculus of admissible inference.

\emph{PROVED} means that an exact implication, identity, certificate, or finite
exhaustion was replayed by the machine.  For example, the four off-grid rows at
\(w\in\{101,103,107,109\}\), \(u=-1\), are PROVED per row after independent
replay; this is a finite horizon statement, not a uniform theorem in \(w\).
% src: RESULTS.md:104; NOTES.md:849-861

\emph{CONDITIONAL} means that every remaining hypothesis is stated and named.
The six-unknown conclusion is conditional on the explicitly isolated global
assembly lemma, so no quantifier record is claimed.  At the analytic end of the
class reduction, the simultaneous primality of \(q_1(t)\) and the degree-\(8\)
polynomial \(F_{\mathrm{cell}}(t)\) awaits Schinzel's Hypothesis H; the finite
local audit does not supply that conjectural prime-value conclusion.
% src: README.md:59-60; README.md:225-230; NOTES.md:862-869

\emph{EVIDENCE} means a measured statistic with its sampling frame attached.
The matched obstruction artifact reports marginal ratio \(0.9993\), a union
residual of \(-2.9/8760\), and a marginal-preserving permutation statistic
\(z=-0.09\) with empirical tails \(0.61/0.53\).  These numbers validate the
matched finite model as EVIDENCE; they do not prove independence or a density
theorem.
% src: data/l17_matched.jsonl:3; NOTES.md:812-817

Labels only move upward through a valid proof.  In particular,
\texttt{PrimalityBound}, \texttt{FactorBudget}, a timeout guard, a
zero-Jacobi outcome, or a missing producer cannot upgrade EVIDENCE to PROVED
and cannot turn an open case into a negative result.  A refusal contributes
\emph{unknown}; only a Fermat witness, a proper factor, or another exact
certificate contributes a mathematical verdict.
% src: README.md:200-204; h10q.py:127-146; NOTES.md:835-845

\section{Proven-primality arithmetic}

\subsection{The deterministic range}
The entry point \texttt{\_is\_prime} first applies strong Miller--Rabin at the
thirteen bases from \(2\) through \(41\)~\cite{jaeschke}.  Its strict deterministic cutoff is
\[
  \texttt{\_MR\_LIMIT}=3317044064679887385961981.
\]
For \(n\) below this bound, passing all bases is a proved primality verdict.
At or above the bound, passing is only a screen and the code proceeds to a
certificate; it never accepts ``probable prime'' as prime.
% src: h10q.py:32-39; h10q.py:169-178

This distinction is load-bearing because factorization feeds exact Hilbert
symbols and ramification sets.  The regression suite therefore includes the
least strong pseudoprime at the cutoff, a prime above the cutoff that must be
Pocklington-certified, hard cases that must refuse, and a balanced semiprime
that must exhaust the factor budget rather than be guessed.
% src: h10q.py:261-297; RESULTS.md:42-53

\subsection{Pocklington and the relaxed threshold}
For a screened \(n\), the engine writes
\[
 n-1=FR,
\]
where every prime dividing the accumulated part \(F\) is recursively
certified.  For each such prime \(q\), it seeks a base \(a\) satisfying
\[
 a^{n-1}\equiv 1\pmod n,
 \qquad
 \gcd\bigl(a^{(n-1)/q}-1,n\bigr)=1.
\]
The generalized Pocklington conclusion is immediate when \(F^2>n\).  A
failed Fermat congruence or a proper gcd proves compositeness; failure to find a
usable witness is a \texttt{PrimalityBound} refusal.
% src: h10q.py:81-91; h10q.py:134-148

Before attempting harder splitting, the implementation trial-divides \(n-1\)
by every prime below \(10^6\).  It computes the integer ceiling of \(n^{1/3}\)
exactly and stops that trial loop as soon as the accumulated certified part
reaches the cube-root threshold.  Uncertifiable recursive chunks and
factor-budget-exhausted chunks are excluded from \(F\), never counted as if
they had factored.
% src: h10q.py:71-78; h10q.py:94-135; README.md:42

At the weaker threshold \(F^3\ge n\), the
Brillhart--Lehmer--Selfridge relaxation, in the form recorded as
Crandall--Pomerance Theorem~\(4.1.5\), applies.  The code first checks the
required coprimality of \(F\) and \(R\), writes \(R=c_2F+c_1\), and tests
\[
 D=c_1^2-4c_2.
\]
Under the certified congruence conditions a composite \(n\) can then have
exactly two factors of the form \(uF+1\) and \(vF+1\).  A nonnegative square
\(D\) with the correct parity reconstructs those two factors and proves
compositeness; otherwise the discriminant test proves primality.  This is a
certificate theorem, not a heuristic strengthening of Miller--Rabin.
% src: h10q.py:85-90; h10q.py:149-166; README.md:42

\subsection{Explicit refusal types}
\texttt{PrimalityBound} means that the screen passed outside its deterministic
range but the certified part or witnesses were insufficient.
\texttt{FactorBudget} means Brent's rho routine exhausted its bounded work on a
large remaining composite.  Callers either propagate these exceptions or mark
the candidate as refused.  The suite contains adversarial paths specifically
to ensure that a partial factorization is not returned as complete and that a
refusal is not counted as a soluble member.
% src: h10q.py:42-48; h10q.py:181-216; h10q.py:281-297

\section{Suite architecture and frozen authority}

A default invocation is \texttt{python3 h10q.py}; the extended invocation adds
\texttt{--extended}.  Both begin with arithmetic self-tests and the ledger
verification, then replay the successive finite layers.  The current engine
note records approximately \(51\) seconds for the default suite and
approximately \(118\) seconds for the extended suite on the recorded run.
% src: h10q.py:4474-4519; README.md:42

The difference is scope, not soundness.  The default L6 replay covers the odd
prime targets below \(100\), while the extended replay reaches the full
\(61\)-row authority below \(300\).  L13 class sampling uses \(12\) escape
rows by default and \(40\) in the extended run; fully soluble escape witnesses
use \(24\) versus all \(60\).  Most importantly, the default L14 check takes a
fixed seeded sample of \(60\) frozen H-class members, while the extended check
replays all \(293\).
% src: h10q.py:1445-1491; h10q.py:4133-4176; h10q.py:4314-4325

\begin{table}[t]
\centering
\small
\caption{Frozen finite authorities and their intended role.}
\begin{tabular}{@{}lll@{}}
\toprule
Authority & Rows & Role \\
\midrule
\texttt{\_L6\_WITNESSES} & \(61\) & bounded canonical tied evidence \\ % src: README.md:43
\texttt{\_L9\_STEERED} & \(345\) & steered cell witnesses, fully replayed \\ % src: README.md:43; RESULTS.md:84
\texttt{\_L11\_CLASSES} & \(190\) & reachable aligned prime classes \\ % src: h10q.py:3329-3344
\texttt{\_L12\_ESCAPE} & \(103\) & non-coprime aligned escape classes \\ % src: h10q.py:3594-3606
\texttt{\_L13\_ESCAPE2} & \(60\) & fully soluble walled-cell witnesses \\ % src: h10q.py:3615-3625
\texttt{\_L13\_RESIDUAL} & \(6\) & legacy pointwise audit slice \\ % src: data/l17_schinzel_audit.jsonl:1; l17_schinzel_audit.py:10-11
\texttt{\_L13\_ZERO\_BAD} & \(16\) & frozen ladder regression rows \\ % src: RESULTS.md:109; h10q.py:4014-4026
\texttt{\_L13\_H\_CLASSES} & \(293\) & one certified member per aligned grid class \\ % src: h10q.py:4011; h10q.py:4314-4318
\bottomrule
\end{tabular}
\end{table}

The first two authorities have a stronger serialization contract.  Each suite
run derives canonical JSONL text from the in-code objects and requires
\path{data/l6_witnesses.jsonl} and
\path{data/l9_steered.jsonl} to match byte for byte.  Drift, absence, or an
underived field fails the run.  The only supported export command is
\begin{verbatim}
python3 h10q.py --export-evidence
\end{verbatim}
which writes those canonical serializations; it does not canonize arbitrary
search logs.
% src: README.md:43; h10q.py:1309-1315; h10q.py:1476-1491; h10q.py:4475-4486

\section{Artifact schemas and licenses}

A schema is part of the proof boundary.  Row types say whether a record is
configuration, a replayable certificate, a failed attempt, or an aggregate.
The principal persisted objects are summarized below; ``licenses'' always
means ``licenses at the recorded finite scope.''

\begin{table}[p]
\centering
\scriptsize
\caption{Principal artifact schemas.}
\begin{tabular}{@{}p{.22\textwidth}p{.31\textwidth}p{.39\textwidth}@{}}
\toprule
Artifact & Row types and key fields & What the row licenses \\
\midrule
\path{l17_sieve.jsonl} &
\texttt{meta}; \texttt{class-summary} with \texttt{cell}, \texttt{family},
\texttt{members}, forbidden residues, rates, alarms; \texttt{global-summary} &
Per-class residue determination and the recorded member window; aggregates are
measured summaries, not tail bounds. \\ % src: data/l17_sieve.jsonl:1-2; data/l17_sieve.jsonl:295

\path{l17_badroots_closures.jsonl} &
\texttt{meta}; \texttt{class} with progression parameters, root lists, exact
or residual masses, exclusions; \texttt{summary} with reconciliation &
Certified bad-signed \(k\bmod p\) roots and their stated \(p\)-adic masses for
the closure progression; excluded roots license only an interval. \\ % src: data/l17_badroots_closures.jsonl:1-2; data/l17_badroots_closures.jsonl:295

\path{l17_matched.jsonl} &
\texttt{meta}, \texttt{meta-methods}, \texttt{global}, \texttt{class}; keys
include matched support, expectations, residuals, window factors, permutation
null, and partial flags &
EVIDENCE for fit on the identical member sample.  It licenses neither
cross-prime independence nor behavior above the recorded prime window. \\ % src: data/l17_matched.jsonl:1-4

\path{l17_classfactors.jsonl} &
Untagged per-class rows: \texttt{cell}, progression parameters,
\texttt{partial}, \texttt{factor\_lower}, \texttt{factor\_upper},
\texttt{factor\_window}, leading bad primes &
Exact bookkeeping for each finite window interval; closure-effort correlation
is descriptive EVIDENCE only. \\ % src: data/l17_classfactors.jsonl:1; NOTES.md:870-874

\path{l17_schinzel_audit.jsonl} &
\texttt{meta}; \texttt{canonical} with the two polynomials, coefficients,
irreducibility certificate, value gcds, local counts; \texttt{residual} with
pointwise witness and \texttt{schinzel\_applicable=false} &
Finite irreducibility, fixed-divisor, pair-gcd, and local-condition checks.
Residual rows are pointwise only; no simultaneous prime values follow. \\ % src: data/l17_schinzel_audit.jsonl:1-2; data/l17_schinzel_audit.jsonl:295-300

\path{l17_horizon*.jsonl} &
Usually \texttt{meta}, \texttt{class-attempt}, \texttt{aligned-class},
\texttt{member}, \texttt{kernel-replay}, an untagged closure row, and
\texttt{summary}; the compact H107 file has only meta and closure &
Only a verified replay plus the exact closure row licenses an off-grid
closure.  Attempt and refusal rows are provenance, not negative evidence. \\ % src: data/l17_horizon101.jsonl:1-446; data/l17_horizon103.jsonl:1-44; data/l17_horizon107.jsonl:1-2; data/l17_horizon109.jsonl:1-390

\path{l13h_all_closures.json} &
Top-level \texttt{n\_cells} and \texttt{records}; each record has
\texttt{family}, \texttt{cell}, \texttt{a}, \texttt{eps}, \texttt{f},
\texttt{q1}, \texttt{N}, \texttt{k\_zero}, \texttt{Q}, \texttt{rung},
\texttt{tied}, \texttt{ramified\_empty}, \texttt{source} &
The canonical closure authority.  Its thirteen-field record licenses the
zero-bad ladder closure; auxiliary tied or ramified fields may themselves be
refused without weakening that license. \\ % src: data/l13h_all_closures.json:1-24; h10q.py:3711-3715
\bottomrule
\end{tabular}
\end{table}

\subsection{Reading rows without overclaiming}
Metadata is a scope declaration, not a certificate by itself.  A
\texttt{meta} row identifies the authority, coordinate convention, engine,
window, and generator assumptions.  The claim-bearing payload is then found
in a class, canonical, kernel-replay, or closure record, while a summary is a
deterministic reduction of those payload rows.  Consequently a reader should
not cite a global total without retaining the class-level artifact from which
it was reduced, and should not cite an attempted class as a certified class.
% src: data/l17_sieve.jsonl:1-2; data/l17_badroots_closures.jsonl:1-2

The canonical closure document makes this lineage especially explicit.  Its
record stores the cell and family, the aligned progression
\(Q=q_1+k_{\mathrm{zero}}N\), the resulting member parameters, the ladder
rung, and the source scan.  The Python authority stores the corresponding
tuple needed for replay.  During L14 verification the suite reconstructs
\(Q\), proves it prime, forms \(b=\varepsilon fQ\), re-runs the smooth
emergent test and cofactor ladder, and accepts only the verdict
\texttt{zero}.  The optional full tied and ramification calculations are
cross-checks rather than hidden prerequisites for this accepted certificate.
% src: data/l13h_all_closures.json:1-24; h10q.py:4321-4368

Large integers and rational parameters are serialized as decimal strings or
explicit numerator--denominator data where the producing schema requires
them.  Status words such as \texttt{bad}, \texttt{deferred},
\texttt{nonprime}, and \texttt{refused:FactorBudget} are intentionally not
coerced to booleans.  This preserves the distinction between a proved
negative result, a skipped member, and an arithmetic refusal.  Schema
compatibility therefore includes both field names and status semantics; a
consumer that collapses those states has changed the estimand even if its JSON
parser succeeds.
% src: data/l17_sieve.jsonl:2; data/l17_horizon109.jsonl:2-5; data/l17_horizon109_crossk.jsonl:2-4

The horizon cross-\(k\) file is deliberately different.  Its meta row freezes
the probe space and determinism contract; untagged rows record primality,
tied, ramified, and size outcomes.  For H109 it contains \(1{,}275\) probes
through \(k\le 50\), but no row with both requested auxiliary successes.  That
is an auxiliary finite non-hit, not a defect in the accepted zero-bad closure.
% src: data/l17_horizon109_crossk.jsonl:1-2; NOTES.md:855-861

\section{Exact reproduction protocol}

Run from the repository directory \texttt{math/h10q}.  The downstream matched
chain is:
\begin{verbatim}
python3 l17_ratemodel_padic.py
python3 l17_badroots_closures.py
python3 l17_matched.py
\end{verbatim}
The middle command imports the checked-in p-adic module and rebuilds closure-
authority roots from the persisted sieve; the final command consumes those
roots and the same persisted member rows.  Thus the first arrow is also a code
dependency.  It does not claim recovery of the original session-side producer
for the older censored-rate artifact.
% src: README.md:209-221; l17_badroots_closures.py:1-27; l17_matched.py:1-24

The independent class-risk and Schinzel audits are regenerated by
\begin{verbatim}
python3 l17_classrisk.py
python3 l17_schinzel_audit.py
\end{verbatim}
The former writes \texttt{data/l17\_classfactors.jsonl} and
\texttt{/tmp/l17\_classrisk.report}; the latter writes
\texttt{data/l17\_schinzel\_audit.jsonl} and its report under
\texttt{/tmp}.
% src: l17_classrisk.py:18-20; l17_classrisk.py:378-387; l17_schinzel_audit.py:33-35; l17_schinzel_audit.py:377-384

The established off-grid horizon is replayed one cell at a time:
\begin{verbatim}
python3 l17_horizon101.py
python3 l17_horizon103.py
python3 l17_horizon107.py
python3 l17_horizon109.py
python3 l17_horizon109_crossk.py
\end{verbatim}
Each closure generator overwrites its named data artifact and emits or writes a
report; the cross-\(k\) generator emits no elapsed or host-dependent fields.
The ledger records about \(17\) minutes for that \(1{,}275\)-probe auxiliary
scan and records byte-identical regeneration for it.  No runtime is supplied
here for commands for which the ledger records none.
% src: README.md:41; NOTES.md:884-893; data/l17_horizon109_crossk.jsonl:1

\section{Review-driven correction as methodology}

The L17 cycle treated review objections as attacks on the estimand, not as
requests for cosmetic caveats.  Five corrections illustrate the method.
% src: NOTES.md:777-883

\paragraph{Recurring-prime-only support.}
The first small-prime product used only primes recurring in the observed
sample.  Omitting nonrecurring bad primes mechanically inflated the predicted
clean mass; it was not evidence of cross-prime correlation.  The fix was to
construct closure-authority root tables across every prime in the declared
window and to target larger primes actually observed in the same cell, with
simple roots assigned exact mass \(1/(p+1)\) and singular roots recursively
resolved or intervalized.
% src: NOTES.md:788-797; l17_badroots_closures.py:1-9

\paragraph{Wrong-sign conditioning.}
A proposed explanation attributed the clean-rate offset to conditioning on
prime \(Q\), but the sign was backwards.  At the \(Q\)-root residue,
\(2\alpha b\equiv0\), so the implementation marks the residue safe rather
than bad.  The corrected interpretation is a difference between Haar measure
on \(\Z_p\) and the finite \(k\)-window.  Its measured size is approximately
\(+0.0195\), or \(3.63\) standard deviations in the artifact; it is
characterized, not promoted to a dependence claim.
% src: NOTES.md:798-807; data/l17_matched.jsonl:3

\paragraph{Cohort versus closure progression.}
The first matcher combined sieve members rebuilt from canonical closure
classes with roots computed on an earlier wave cohort.  In the documented ESC
example at cell \((3,(-3,1))\), the cohort used \(q_1=41\), \(\varepsilon=+1\),
whereas the closure used \(q_1=59\), \(\varepsilon=-1\).  The fix was to
rebuild every root in the closure record's own \((\varepsilon,f,q_1,N)\)
coordinate.  The repaired table reconciles \(5{,}375/5{,}375\) observed
odd-negative occurrences.
% src: NOTES.md:818-830; data/l17_badroots_closures.jsonl:295

\paragraph{Discarded resolved mass.}
An excluded singular root can have a rigorously resolved prefix even when its
remaining Taylor tree is unfinished.  An intermediate matcher discarded that
known mass along with the unresolved remainder.  The corrected artifact stores
both pieces and forms the interval
\[
 \left[\prod_p(1-k_p-r_p),\ \prod_p(1-k_p)\right].
\]
Across the \(67\) excluded pairs the maximum unresolved residual is
\(6.9\times10^{-15}\), so the finite-window interval is valid without
pretending to control any larger prime.
% src: NOTES.md:875-883; data/l17_matched.jsonl:1

\paragraph{The vacuous integer-support p-value.}
A clipped two-sided empirical value on discrete integer support reported
\(1.0\), which carried little diagnostic information.  The final report keeps
the marginal-preserving null but presents its standardized statistic and both
empirical tails.  Together with the repaired support and masses, the final
matched values are marginal ratio \(0.9993\), union residual \(-2.9/8760\),
\(z=-0.09\), and tails \(0.61/0.53\).  These are the final EVIDENCE numbers;
the abandoned instruments license nothing.
% src: NOTES.md:879-883; data/l17_matched.jsonl:3; l17_matched.py:162-166

\section{Limitations}

\begin{enumerate}
\item Five older L17 files are persisted evidence without checked-in original
producers: \path{l17_sieve.jsonl}, the cohort
\path{l17_badroots.jsonl}, \path{l17_ratemodel_censored.jsonl},
\path{l17_cofactor.jsonl}, and
\path{l17_cofactor_nonclosure.jsonl}.  Downstream scripts consume reviewed
copies, but the original session-side acquisition scripts were not recovered.
% src: README.md:207-221; NOTES.md:884-894

\item The class factors and matched Haar comparison are statistics only for
\(p\le10^4\).  They provide no bound whatsoever for primes above \(10^4\), so
they cannot by themselves prove a positive infinite product or a uniform
density.
% src: README.md:35; data/l17_matched.jsonl:1

\item The nonclosure cofactor distribution is measurement-limited.  Only
\(11/400\) sampled rows were exactly resolved and \(389\) were prover
refusals; the predeclared power gate failed.  The distributional conclusion is
therefore INCONCLUSIVE, even though parity remained consistent on resolved
rows.
% src: README.md:38; NOTES.md:835-845

\item Some closure records contain auxiliary tied-status or ramification
refusals.  For example, the H109 closure records budget/refusal values in those
fields while its proved zero-bad ladder verdict remains the accepted closure
form.  Such fields must not be silently rewritten as success, and their
refusals are not supporting evidence.
% src: RESULTS.md:104; data/l17_horizon109.jsonl:387-390

\item The Schinzel audit proves finite local admissibility and irreducibility
for the canonical pairs, not simultaneous prime values of the degree-\(8\)
polynomials.  The record claim remains CONDITIONAL, and the finite off-grid
horizon does not imply the all-\(w\) hypothesis.
% src: README.md:225-231; NOTES.md:862-869
\end{enumerate}

The reproducibility claim is therefore intentionally narrow but useful:
checked-in authorities can be replayed, specified downstream artifacts can be
regenerated, and every refusal remains visible.  The paper's mathematical
claims are exactly those finite certificates labeled PROVED; its statistics
remain EVIDENCE; and its analytic extrapolations remain CONDITIONAL.
% src: README.md:12-43; README.md:207-221

\begin{thebibliography}{9}

\bibitem{cp}
R.~Crandall and C.~Pomerance,
\emph{Prime Numbers: A Computational Perspective}, 2nd ed.,
Springer, New York, 2005.
% src: h10q.py:82-90 (the in-repository citation is CP, Theorems 4.1.3 and 4.1.5)

\bibitem{jaeschke}
G.~Jaeschke,
\emph{On strong pseudoprimes to several bases},
Mathematics of Computation \textbf{61} (1993), no.~204, 915--926.
% src: h10q.py deterministic Miller-Rabin base sets below _MR_LIMIT

\end{thebibliography}

\end{document}
